\textbf{1--4. Boundary conditions.} The fixed end has zero displacement and slope. The free end has zero moment and shear. With the B.1 harmonic solution, these conditions become
\[
W(0)=W'(0)=W''(L)=W'''(L)=0,\quad
W=C_1\sin\gamma x+C_2\cos\gamma x+C_3\sinh\gamma x+C_4\cosh\gamma x.
\]
The harmonic displacement is $w(x,t)=W(x)e^{-i\omega t}$. The conditions at $x=0$ require $C_2+C_4=0$ and $\gamma(C_1+C_3)=0$. For a flexible mode $\gamma>0$, so $C_3=-C_1$ and $C_4=-C_2$. The remaining shape is
\[
W=C_1(\sin\gamma x-\sinh\gamma x)+C_2(\cos\gamma x-\cosh\gamma x).
\]
Writing $z=\gamma L$, the zero-moment and zero-shear conditions at the free end then give
\[
\begin{bmatrix}
\sin z+\sinh z & \cos z+\cosh z\\
\cos z+\cosh z & -\sin z+\sinh z
\end{bmatrix}
\begin{bmatrix}C_1\\C_2\end{bmatrix}=\mathbf0.
\]

\textbf{5--8. Eigenvalue equation.} $C_1=C_2=0$ produces the trivial zero-displacement solution. A nonzero mode requires a zero determinant. Expanding the determinant and applying $\sin^2z+\cos^2z=1$ and $\cosh^2z-\sinh^2z=1$ gives
\[
(\sin z+\sinh z)(-\sin z+\sinh z)-(\cos z+\cosh z)^2
=-2(1+\cos z\cosh z)=0.
\]
The characteristic equation is therefore
\[
\boxed{\cos z\cosh z+1=0.}
\]
For numerical root finding I use the equivalent, well-scaled function $D(z)=\cos z+1/\cosh z$. The first five positive roots are
\[
\boxed{z_1=1.87510407,\quad z_2=4.69409113,\quad z_3=7.85475744,\quad z_4=10.99554073,\quad z_5=14.13716839.}
\]

\textbf{9--13. Frequencies.} Since $\gamma=\sqrt\omega/a$, the wavenumbers and natural frequencies are
\[
\boxed{\gamma_j=\frac{z_j}{L},\quad
\omega_j=\frac{a^2z_j^2}{L^2}
=z_j^2\sqrt{\frac{EI}{mL^4}},\quad
f_j=\frac{z_j^2}{2\pi}\sqrt{\frac{EI}{mL^4}}.}
\]

\textbf{14--19. Modes and orthonormality.} Either free-end equation yields
\[
\frac{C_1}{C_2}=-\beta_j,\qquad
\beta_j=\frac{\cosh z_j+\cos z_j}{\sinh z_j+\sin z_j}
=\frac{\sinh z_j-\sin z_j}{\cosh z_j+\cos z_j}.
\]
Thus $C_1=-\beta_j C_2$, $C_3=\beta_j C_2$, and $C_4=-C_2$. Setting the requested $C_2=1$ gives
\[
\boxed{W_j(x)=\cos\gamma_jx-\cosh\gamma_jx
+\beta_j(\sinh\gamma_jx-\sin\gamma_jx).}
\]
The course notes use the opposite overall sign. Both choices give the same nodes and natural frequencies. Integrating the two mode equations against one another and subtracting gives
\[
(\gamma_j^4-\gamma_k^4)\int_0^L W_jW_k\,dx
=\left[W_j'''W_k-W_j''W_k'+W_j'W_k''-W_jW_k'''\right]_0^L.
\]
The clamp sets both displacement and slope to zero at the left end, and the free end sets both curvature and its derivative to zero at the right. The boundary term therefore vanishes. Distinct modes have zero overlap, while the course-note scale above gives $\int_0^L W_j^2\,dx=L$. Hence $L^{-1}\int_0^L W_jW_k\,dx=\delta_{jk}$. Multiplying by the uniform mass per unit length gives the modal mass $m_j=mL$ for this scale.

\textbf{20. High-order roots.} For $j\ge6$, $z_j\simeq(j-\tfrac12)\pi$. The root calculation gives the following comparison. The difference is numerical root minus approximation.
\begin{center}
\begin{tabular}{rrrr}
\toprule
$j$ & Numerical $z_j$ & Approximation & Difference\\
\midrule
1 & 1.87510407 & 1.57079633 & $3.04\times10^{-1}$\\
2 & 4.69409113 & 4.71238898 & $-1.83\times10^{-2}$\\
3 & 7.85475744 & 7.85398163 & $7.76\times10^{-4}$\\
4 & 10.99554073 & 10.99557429 & $-3.36\times10^{-5}$\\
5 & 14.13716839 & 14.13716694 & $1.45\times10^{-6}$\\
6 & 17.27875953 & 17.27875959 & $-6.27\times10^{-8}$\\
7 & 20.42035225 & 20.42035225 & $2.71\times10^{-9}$\\
8 & 23.56194490 & 23.56194490 & $-1.17\times10^{-10}$\\
9 & 26.70353756 & 26.70353756 & $5.06\times10^{-12}$\\
10 & 29.84513021 & 29.84513021 & $-2.20\times10^{-13}$\\
\bottomrule
\end{tabular}
\end{center}
The alternating difference decreases rapidly, consistent with the asymptotic approximation in the course notes, PDF pages 179--184.

\lstinputlisting[style=hwmatlab,caption={Cantilever eigenvalues and high-order approximation.}]{matlab/B2.m}
